% ══════════════════════════════════════════════════════════════
%  Modular Period Structure, Fractal Laws, and Algebraic
%  Universality of the Affine Sequence  T_k = 6·2^k + 1
%
%  Paket 2 — Number Theory Preprint (v3.0 — Full 28-module edition)
%  Author: Ömer Çetintaş
%  Date:   March 2026
% ══════════════════════════════════════════════════════════════
\documentclass[11pt, a4paper]{article}

% ── Encoding & Fonts ──────────────────────────────────────────
\usepackage[utf8]{inputenc}
\usepackage[T1]{fontenc}
\usepackage{lmodern}
\usepackage{microtype}

% ── Page Geometry ─────────────────────────────────────────────
\usepackage[
  a4paper,
  top=2.5cm,
  bottom=2.5cm,
  left=2.5cm,
  right=2.5cm,
  headheight=14pt
]{geometry}

% ── Mathematics ───────────────────────────────────────────────
\usepackage{amsmath,amssymb,amsthm,mathtools}

% ── Colors ────────────────────────────────────────────────────
\usepackage[dvipsnames]{xcolor}

% ── Tables & Floats ──────────────────────────────────────────
\usepackage{booktabs}
\usepackage{array}
\usepackage{float}
\usepackage[font=small, labelfont=bf, labelsep=period]{caption}

% ── Header / Footer ──────────────────────────────────────────
\usepackage{fancyhdr}
\pagestyle{fancy}
\fancyhf{}
\fancyhead[L]{\small\textit{Modular Periods of $T_k = 6 \cdot 2^k + 1$}}
\fancyhead[R]{\small\textit{\c{C}etinta\c{s}, 2026}}
\fancyfoot[C]{\small\thepage}
\renewcommand{\headrulewidth}{0.4pt}
\renewcommand{\footrulewidth}{0pt}

\fancypagestyle{firstpage}{%
  \fancyhf{}
  \fancyfoot[C]{\small\thepage}
  \renewcommand{\headrulewidth}{0pt}
}

% ── Hyperlinks ────────────────────────────────────────────────
\usepackage[
  colorlinks=true,
  linkcolor=blue!70!black,
  citecolor=blue!70!black,
  urlcolor=blue!70!black,
  pdfauthor={\"Omer \c{C}etinta\c{s}},
  pdftitle={Modular Period Structure and Algebraic Universality of T_k = 6*2^k + 1},
  pdfsubject={Number Theory, Formal Verification},
  pdfkeywords={OMARWA sequence, multiplicative order, Lean 4, modular periodicity}
]{hyperref}

% ── Theorem Environments ─────────────────────────────────────
\theoremstyle{plain}
\newtheorem{theorem}{Theorem}[section]
\newtheorem{lemma}[theorem]{Lemma}
\newtheorem{proposition}[theorem]{Proposition}
\newtheorem{corollary}[theorem]{Corollary}

\theoremstyle{definition}
\newtheorem{definition}[theorem]{Definition}
\newtheorem{conjecture}[theorem]{Conjecture}

\theoremstyle{remark}
\newtheorem{remark}{Remark}[section]

% ── Custom Commands ───────────────────────────────────────────
\newcommand{\ord}{\operatorname{ord}}
\DeclareMathOperator{\lcm}{lcm}
\DeclareMathOperator{\DR}{DR}

% ── Spacing ───────────────────────────────────────────────────
\usepackage{enumitem}
\setlist[enumerate]{leftmargin=2em, itemsep=0.3em}
\setlist[itemize]{leftmargin=2em, itemsep=0.2em}

% ══════════════════════════════════════════════════════════════
\begin{document}
\thispagestyle{firstpage}

% ── TITLE BLOCK ───────────────────────────────────────────────
\begin{center}
  \vspace*{0.3cm}

  {\LARGE\bfseries
    Modular Period Structure, Fractal Laws,\\[0.3cm]
    and Algebraic Universality of the\\[0.3cm]
    Affine Sequence $T_k = 6 \cdot 2^k + 1$}

  \vspace{0.8cm}

  {\large \"Omer \c{C}etinta\c{s}}\\[0.3cm]
  {\small ORCID: \href{https://orcid.org/0009-0003-2977-6048}{0009-0003-2977-6048}}\\[0.15cm]
  {\small\textit{Independent Researcher}}

  \vspace{0.6cm}
  {\small March 2026}
\end{center}

\vspace{0.6cm}

% ── ABSTRACT ──────────────────────────────────────────────────
\begin{abstract}
We study the integer sequence $T_k = 6 \cdot 2^k + 1$ (OEIS A004119)
and establish a complete theory of its modular periodicity.
For each positive integer~$m$, we define the OMARWA period $P(m)$ as the
minimal period of the sequence $T_k \bmod m$ and prove the
\textbf{Core Reduction Theorem}: $P(m) = \ord_{m'}(2)$ where
$m' = m / \gcd(m, 6)$.
We derive \textbf{fractal period laws}
$P(p^n) = \ord_p(2) \cdot p^{n-1}$ for non-Wieferich primes $p \nmid 6$
and prove a \textbf{Product Theorem} via the Chinese Remainder Theorem.
These results yield a \textbf{super-period} $L = 30$, derived through
three independent paths, coinciding with the Coxeter number $h(E_8)$.

We further establish a \textbf{Forbidden Unity Theorem} ($T_k \not\equiv 1
\pmod{m}$ when $\gcd(m,6) = 1$), a \textbf{CRT intertwining map} for
the modulus $354 = 2 \times 3 \times 59$, \textbf{Mersenne--Fermat
period dichotomies}, a \textbf{palindromic involution} on residue
sequences, \textbf{network topology identities}, and connections to
\textbf{Weyl exponents of~$E_8$}. The Coxeter numbers of all five
exceptional Lie groups arise as specific OMARWA periods or derived
invariants (\textbf{Exceptional Universality}). We discuss topological
interpretations including \textbf{winding numbers on~$T^3$},
\textbf{Chern--Simons level}, and \textbf{Hopf fibration} structures.
All principal theorems are formally verified in Lean~4 with Mathlib
(956~theorems, zero \texttt{sorry} axioms, 28~modules;
source: \href{https://github.com/omarwa622/OmarwaCore}{github.com/omarwa622/OmarwaCore}).
\end{abstract}

\smallskip
\noindent\textbf{Keywords:} integer sequence $\cdot$ multiplicative
order $\cdot$ modular periodicity $\cdot$ super-period $\cdot$
fractal period law $\cdot$ Coxeter number $\cdot$
affine recurrence $\cdot$ Lean~4 $\cdot$ formal verification

\smallskip
\noindent\textbf{MSC 2020:} 11B50, 11A07, 11A25, 11Y55, 68V20

\vspace{0.5cm}

% ══════════════════════════════════════════════════════════════
\section{Introduction}
% ══════════════════════════════════════════════════════════════

The sequence $T_k = 6 \cdot 2^k + 1$ for $k \geq 0$ produces the terms
$7, 13, 25, 49, 97, 193, 385, \ldots$ (OEIS A004119, shifted:
$T_k = a(k{+}2)$ where $a(n) = 3 \cdot 2^{n-1} + 1$). Despite its
elementary definition as an affine transformation of powers of~$2$, the
modular arithmetic of this sequence exhibits a remarkably rich and
structured periodicity theory.

In this paper we establish the following principal results:

\begin{enumerate}
  \item \textbf{Core Reduction} (Theorem~\ref{thm:core}): $P(m) =
    \ord_{m'}(2)$ where $m' = m/\gcd(m,6)$.

  \item \textbf{Fractal Period Laws}
    (Theorems~\ref{thm:ternary}--\ref{thm:general-fractal}):
    $P(p^n) = \ord_p(2) \cdot p^{n-1}$ for non-Wieferich $p \nmid 6$.

  \item \textbf{Product Theorem} (Theorem~\ref{thm:product}):
    $P(m_1 m_2) = \lcm(P(m_1), P(m_2))$ for $\gcd(m_1, m_2) = 1$.

  \item \textbf{Super-Period} $L = 30 = h(E_8)$
    (Theorems~\ref{thm:sp1}--\ref{thm:sp3}).

  \item \textbf{Sabbath Node} (Theorem~\ref{thm:sabbath}):
    $T_6 = 385 = 5 \times 7 \times 11$, the triple-reset point.

  \item \textbf{Mersenne--Fermat Periods}
    (Theorems~\ref{thm:mersenne}--\ref{thm:fermat}):
    $P(2^p - 1) = p$ for Mersenne primes; $P(17) = 8 = \mathrm{rank}(E_8)$.

  \item \textbf{Forbidden Unity} (Theorem~\ref{thm:forbidden}):
    $T_k \not\equiv 1 \pmod{m}$ when $\gcd(m,6) = 1$.

  \item \textbf{CRT Intertwining} (Section~\ref{sec:crt}):
    The 354-anatomy and the intertwining map
    $\iota(s) = (60s + 295) \bmod 354$.

  \item \textbf{Palindromic Involution, Network Topology, and
    Holon Crystallography}
    (Sections~\ref{sec:palindrome}--\ref{sec:holon}).

  \item \textbf{Weyl Exponents and Exceptional Universality}
    (Sections~\ref{sec:weyl}--\ref{sec:exceptional}).

  \item \textbf{Topological Interpretations}
    (Section~\ref{sec:topology}): winding numbers, Chern--Simons level,
    Hopf fibration.
\end{enumerate}

All results are accompanied by formal Lean~4 proofs
(Appendix~\ref{app:lean}) and computational verification
(Appendix~\ref{app:comp}).

The paper is organized as follows.
Section~2 defines the sequence and establishes basic properties
including the Sabbath node.
Section~3 develops the fractal period theory, Mersenne--Fermat
dichotomy, forbidden unity theorem, and super-period.
Section~4 treats the affine operator and $6n{+}1$ skeleton.
Section~5 presents the Product Theorem and CRT intertwining.
Section~6 examines centered hexagonal intersections.
Section~7 presents the digital root structure.
Section~8 discusses Pascal--Sierpi\'nski--Fibonacci connections and the
Collatz bridge.
Section~9 develops palindromic symmetry and fractal holon structures.
Section~10 establishes network topology and holon crystallography.
Section~11 treats Weyl exponents and exceptional universality.
Section~12 discusses topological and physical interpretations.
Section~13 states open questions and conclusions.

% ══════════════════════════════════════════════════════════════
\section{Definition and Basic Properties of \texorpdfstring{$T_k$}{Tk}}
% ══════════════════════════════════════════════════════════════

\subsection{Definition}

\begin{definition}[OMARWA Sequence]\label{def:omarwa}
The \emph{OMARWA sequence} is defined by $T_k = 6 \cdot 2^k + 1$ for
$k \geq 0$.
\end{definition}

\begin{table}[H]
\centering
\caption{First twelve terms of the OMARWA sequence.}\label{tab:terms}
\smallskip
\begin{tabular}{@{}l*{12}{c}@{}}
\toprule
$k$   & 0 & 1  & 2  & 3  & 4  & 5   & 6   & 7   & 8    & 9    & 10   & 11   \\
\midrule
$T_k$ & 7 & 13 & 25 & 49 & 97 & 193 & 385 & 769 & 1537 & 3073 & 6145 & 12289 \\
\bottomrule
\end{tabular}
\end{table}

\begin{theorem}[Affine Recurrence]\label{thm:recurrence}
$T_{k+1} = 2T_k - 1$ for all $k \geq 0$, with $T_0 = 7$.
\end{theorem}

\begin{proof}
$2T_k - 1 = 2(6 \cdot 2^k + 1) - 1 = 6 \cdot 2^{k+1} + 1 = T_{k+1}$.
\end{proof}

\begin{theorem}[Proth Form]\label{thm:proth}
$T_k = 3 \cdot 2^{k+1} + 1$; since $3 < 2^{k+1}$ for $k \geq 1$, each
such $T_k$ is a Proth number~\cite{proth}.
\end{theorem}

\begin{theorem}[Closed-Form Inverse]\label{thm:inverse}
$k = \log_2\!\bigl(\frac{T_k - 1}{6}\bigr)$ for
$T_k \in \{T_k\}_{k \geq 0}$.
\end{theorem}

\subsection{Modular Periods}

\begin{definition}[OMARWA Period]\label{def:period}
For $m \geq 1$, the \emph{OMARWA period} $P(m)$ is the minimal positive
integer~$p$ such that $T_{k+p} \equiv T_k \pmod{m}$ for all $k \geq 0$.
\end{definition}

\begin{theorem}[Mod 3 Triviality]\label{thm:mod3}
$T_k \equiv 1 \pmod{3}$ for all~$k$. Hence $P(3) = 1$.
\end{theorem}

\begin{proof}
$6 \cdot 2^k \equiv 0 \pmod{3}$, so $T_k = 6 \cdot 2^k + 1 \equiv 1$.
\end{proof}

\begin{theorem}[Full Residue System mod 11]\label{thm:fullres}
$P(11) = 10 = \varphi(11)$, and the residues $T_k \bmod 11$ for
$0 \leq k < 10$ traverse $\mathbb{Z}_{11} \setminus \{1\}$.
\end{theorem}

\begin{proof}
Since $\ord_{11}(2) = 10 = \varphi(11)$, the element~$2$ is a
primitive root modulo~$11$. The map $k \mapsto 6 \cdot 2^k + 1 \bmod 11$
hits every residue except~$1$. Explicitly:
$\{T_0, \ldots, T_9\} \bmod 11 = \{7, 2, 3, 5, 9, 6, 0, 10, 8, 4\}$.
\end{proof}

\subsection{The Sabbath Node}\label{sec:sabbath}

\begin{theorem}[Sabbath Node]\label{thm:sabbath}
$T_6 = 385 = 5 \times 7 \times 11$. In particular, $T_6$ is
simultaneously divisible by~$5$, $7$, and~$11$:
\[
  T_6 \equiv 0 \pmod{5}, \quad
  T_6 \equiv 0 \pmod{7}, \quad
  T_6 \equiv 0 \pmod{11}.
\]
\end{theorem}

\begin{proof}
$T_6 = 6 \cdot 2^6 + 1 = 385$. The factorization $385 = 5 \times 7
\times 11$ is verified directly.
\end{proof}

The triple-reset at $k = 6$ (the seventh term, hence ``Sabbath node'')
is explained by the confluence of periods: $P(5) = 4 \mid 4$, but the
mod~$7$ and mod~$11$ orbits each pass through zero at $k = 6$
(since $P(7) = 3 \mid 6$ and $6 < P(11) = 10$ with
$T_6 \bmod 11 = 0$).

\begin{proposition}[Primality Map]\label{prop:primality}
Among $T_0, T_1, \ldots, T_{29}$, the primes are exactly
$T_k$ for $k \in \{0, 1, 4, 5, 7, 11, 17, 29\}$, yielding eight
primes within the first super-period.
\end{proposition}

\begin{proposition}[Emirp Pair]\label{prop:emirp}
The integers $37$ and $73$ form an emirp pair (both prime, each the
digit-reversal of the other). Both appear naturally:
$37 = P(37) = 36$ gives $\ord_{37}(2) = 36 = \varphi(37)$, and
$73 = T_3 + T_1 + T_0 - 2$.
\end{proposition}

% ══════════════════════════════════════════════════════════════
\section{Fractal Period Laws, Forbidden Unity, and the Super-Period}
% ══════════════════════════════════════════════════════════════

\subsection{Prime-Power Periods}

\begin{theorem}[Ternary Fractal Law]\label{thm:ternary}
For $n \geq 2$: $P(3^n) = 2 \cdot 3^{n-2}$.
\end{theorem}

\begin{proof}
By the Core Reduction (Theorem~\ref{thm:core}),
$P(3^n) = \ord_{3^{n-1}}(2)$. Since $\ord_3(2) = 2$ and $2^2 = 4
\equiv 1 \pmod{3}$ but $4 \not\equiv 1 \pmod{9}$, the
lifting-the-exponent result~\cite{narkiewicz} gives
$\ord_{3^{n-1}}(2) = 2 \cdot 3^{n-2}$.
\end{proof}

\begin{theorem}[Undecimal Fractal Law]\label{thm:undecimal}
For $n \geq 1$: $P(11^n) = 10 \cdot 11^{n-1}$.
\end{theorem}

\begin{theorem}[General Fractal Law]\label{thm:general-fractal}
Let $p \nmid 6$ be a prime satisfying the non-Wieferich condition
$2^{\ord_p(2)} \not\equiv 1 \pmod{p^2}$. Then for $n \geq 1$:
\[
  P(p^n) = \ord_p(2) \cdot p^{n-1}.
\]
\end{theorem}

\begin{proof}
Since $p \nmid 6$, the Core Reduction gives
$P(p^n) = \ord_{p^n}(2)$. Under the non-Wieferich hypothesis, the
standard lifting result~\cite{narkiewicz} gives
$\ord_{p^n}(2) = \ord_p(2) \cdot p^{n-1}$.
\end{proof}

\begin{remark}\label{rem:wieferich}
The only known Wieferich primes are $1093$ and $3511$~\cite{wieferich}.
All other primes $p < 10^{17}$ satisfy the hypothesis.
\end{remark}

\subsection{Core Reduction}

\begin{theorem}[Core Reduction]\label{thm:core}
For all $m \geq 1$ with $\gcd(2,\, m/\gcd(m,6)) = 1$:
\[
  P(m) = \ord_{m'}(2), \qquad m' = m / \gcd(m, 6).
\]
\end{theorem}

\begin{proof}
Write $g = \gcd(m, 6)$. The period of $T_k = 6 \cdot 2^k + 1$
modulo~$m$ equals the period of $6 \cdot 2^k$ modulo~$m$, since adding
a constant does not change the period. Factoring:
$6 \cdot 2^k \bmod m$ has the same period as $2^k \bmod m'$ because
$(6/g) \cdot 2^k \bmod (m/g)$---with $(6/g)$ invertible
modulo~$m/g$---inherits the period of~$2^k$.
\end{proof}

\subsection{Mersenne and Fermat Periods}\label{sec:mersenne}

The OMARWA periods at Mersenne and Fermat primes exhibit a clean
dichotomy:

\begin{theorem}[Mersenne Periods]\label{thm:mersenne}
For Mersenne primes $M_p = 2^p - 1$:
\[
  P(M_p) = p.
\]
In particular: $P(7) = 3$, $P(31) = 5$, $P(127) = 7$.
\end{theorem}

\begin{proof}
By the Core Reduction, $P(M_p) = \ord_{M_p}(2)$. Since
$2^p \equiv 1 \pmod{M_p}$ by definition and $M_p$ is prime,
$\ord_{M_p}(2) \mid p$. Since $p$ is prime and $\ord_{M_p}(2) > 1$,
the only possibility is $\ord_{M_p}(2) = p$.
\end{proof}

\begin{theorem}[Fermat--Rank Lock]\label{thm:fermat}
$P(17) = 8 = \mathrm{rank}(E_8)$. More generally, for the Fermat prime
$F_n = 2^{2^n} + 1$, we have $\ord_{F_n}(2) = 2^{n+1}$.
\end{theorem}

\begin{proof}
$\ord_{17}(2) = 8$ since $2^8 = 256 = 15 \times 17 + 1 \equiv 1
\pmod{17}$ and $2^4 = 16 \equiv -1 \not\equiv 1 \pmod{17}$.
\end{proof}

The coincidence $P(17) = 8 = \mathrm{rank}(E_8)$ is the
``Rank Lock'' connecting the OMARWA period at the Fermat prime~$17$
to the rank of the exceptional Lie algebra~$E_8$.

\subsection{Forbidden Unity Theorem}\label{sec:forbidden}

\begin{theorem}[Forbidden Unity]\label{thm:forbidden}
For any $m \geq 2$ with $\gcd(m, 6) = 1$:
\[
  T_k \not\equiv 1 \pmod{m} \quad \text{for all } k \geq 0.
\]
\end{theorem}

\begin{proof}
Suppose $T_k \equiv 1 \pmod{m}$. Then $6 \cdot 2^k \equiv 0 \pmod{m}$.
Since $\gcd(m, 6) = 1$ and $\gcd(m, 2^k) = 1$ (as $m$ is odd and
coprime to~$3$), we have $m \mid 6 \cdot 2^k$ impossible for
$m \geq 2$. Contradiction.
\end{proof}

\begin{corollary}
The residue $1$ is \emph{forbidden} in the mod~$11$ orbit of $T_k$:
$\{T_k \bmod 11 : k \geq 0\} = \mathbb{Z}_{11} \setminus \{1\}$.
Similarly for $m = 29, 59, 177$.
\end{corollary}

\subsection{Lifted Super-Period}

\begin{proposition}[Lifted Super-Period]\label{prop:lifted}
$P(121) = P(11^2) = 110$ and
$L' = \lcm(P(9), P(121), P(27)) = \lcm(2, 110, 6) = 330 = 11 \times
h(E_8)$.
\end{proposition}

\subsection{The Super-Period}\label{sec:super}

\begin{definition}[Super-Period]\label{def:superperiod}
The \emph{super-period} $L$ is the minimal positive integer such that
$T_{k+L} \equiv T_k$ simultaneously modulo~$9$, $11$, and~$27$.
\end{definition}

\begin{theorem}[Super-Period, Path~I]\label{thm:sp1}
$L = \lcm(P(9), P(11), P(27)) = \lcm(2, 10, 6) = 30.$
\end{theorem}

\begin{theorem}[Super-Period, Path~II]\label{thm:sp2}
$P(77) = P(7 \times 11) = \lcm(P(7), P(11)) = \lcm(3, 10) = 30.$
\end{theorem}

\begin{theorem}[Super-Period, Path~III]\label{thm:sp3}
$P(99) = P(9 \times 11) = \lcm(P(9), P(11)) = \lcm(2, 10) = 10.$ \\
Yet $\lcm(P(99), P(27)) = \lcm(10, 6) = 30 = L$.
\end{theorem}

\subsection{Composite Modular Periods}

\begin{proposition}\label{prop:composite}
$P(33) = 10$, \; $P(99) = 10$, \; $P(407) = 180$, \;
$P(777) = 36$, \; $P(999) = 36$, \; $P(1453) = 1452$.
\end{proposition}

The value $P(1453) = 1452 = \varphi(1453)$ means $2$ is a
primitive root modulo~$1453$, consistent with Artin's
conjecture~\cite{artin}.

% ══════════════════════════════════════════════════════════════
\section{Affine Operator and the \texorpdfstring{$6n{+}1$}{6n+1}
  Skeleton}\label{sec:affine}
% ══════════════════════════════════════════════════════════════

\begin{definition}[$6n{+}1$ Skeleton]\label{def:skeleton}
Define $A_n = 6n + 1$ for $n \geq 0$. Then:
\end{definition}

\begin{theorem}[$6n{+}1$ Confinement]\label{thm:confinement}
$T_k \equiv 1 \pmod{6}$ for all $k \geq 0$.
\end{theorem}

\begin{proof}
$T_k = 6 \cdot 2^k + 1 \equiv 0 + 1 = 1 \pmod{6}$.
\end{proof}

\begin{theorem}[Sampling Theorem]\label{thm:sampling}
$T_k = A_{2^k}$; the OMARWA sequence samples the $6n{+}1$ skeleton at
powers of~$2$.
\end{theorem}

\begin{proof}
$A_{2^k} = 6 \cdot 2^k + 1 = T_k$.
\end{proof}

For each modulus $m$ with $m' = m/\gcd(m,6)$, the doubling map
$\phi : x \mapsto 2x \pmod{m'}$ acts on $\mathbb{Z}_{m'}$.

\begin{proposition}[Orbit-Period Identity]\label{prop:orbit}
The length of the orbit of~$1$ under $\phi$ in
$\mathbb{Z}_{m'}$ equals $P(m)$.
\end{proposition}

\begin{proposition}[Key Periods from the Affine Operator]\label{prop:key-periods}
\leavevmode
\begin{itemize}
  \item $P(7) = 3 = \ord_7(2)$,
  \item $P(11) = 10 = \ord_{11}(2)$,
  \item $P(17) = 8 = \ord_{17}(2) = \mathrm{rank}(E_8)$,
  \item $P(29) = 28$ (half-ribbon period),
  \item $P(59) = 58$ (full ribbon traversal),
  \item $P(77) = 30 = h(E_8)$ (Coxeter Lock).
\end{itemize}
\end{proposition}

% ══════════════════════════════════════════════════════════════
\section{Product Theorem and CRT Intertwining}\label{sec:crt}
% ══════════════════════════════════════════════════════════════

\subsection{Product Theorem}

\begin{theorem}[Product Theorem / CRT]\label{thm:product}
For $\gcd(m_1, m_2) = 1$:
\[
  P(m_1 m_2) = \lcm(P(m_1), P(m_2)).
\]
\end{theorem}

\begin{proof}
By the Chinese Remainder Theorem~\cite{ireland}, the ring
$\mathbb{Z}_{m_1 m_2} \cong \mathbb{Z}_{m_1} \times \mathbb{Z}_{m_2}$,
and the period of a pair $(a_k \bmod m_1,\; a_k \bmod m_2)$ is the
least common multiple of the component periods.
\end{proof}

\begin{theorem}[GCD Law]\label{thm:gcd}
For $i \neq j$: $\gcd(T_i, T_j) \mid 2^{|i-j|} - 1$.
\end{theorem}

\begin{proof}
$T_i - T_j = 6 \cdot 2^{\min(i,j)}(2^{|i-j|} - 1)$.
If $d = \gcd(T_i, T_j)$, then $\gcd(d, 6) = 1$ (since $T_k$ is odd
and $T_k \not\equiv 0 \pmod{3}$) and $\gcd(d, 2^{\min(i,j)}) = 1$.
Hence $d \mid (2^{|i-j|} - 1)$.
\end{proof}

\begin{corollary}\label{cor:coprime}
$\gcd(T_k, T_{k+1}) = 1$ for all $k \geq 0$.
\end{corollary}

\begin{table}[H]
\centering
\caption{Selected values of the OMARWA period $P(m)$.}\label{tab:periods}
\smallskip
\begin{tabular}{@{}rrcrl@{}}
\toprule
$m$ & $P(m)$ & $m'$ & $\ord_{m'}(2)$ & Note \\
\midrule
3  & 1  & 1  & ---  & $T_k \equiv 1$ \\
5  & 4  & 5  & 4    & --- \\
7  & 3  & 7  & 3    & Mersenne $2^3 - 1$ \\
9  & 2  & 3  & 2    & Binary alternation \\
11 & 10 & 11 & 10   & Primitive root \\
13 & 12 & 13 & 12   & --- \\
17 & 8  & 17 & 8    & Fermat; $= \mathrm{rank}(E_8)$ \\
27 & 6  & 9  & 6    & $= h(G_2)$ \\
29 & 28 & 29 & 28   & Half-ribbon \\
31 & 5  & 31 & 5    & Mersenne $2^5 - 1$ \\
37 & 36 & 37 & 36   & Primitive root \\
59 & 58 & 59 & 58   & Ribbon width \\
77 & 30 & 77 & 30   & $= h(E_8)$ \\
81 & 18 & 27 & 18   & $= h(E_7)$ \\
121 & 110 & 121 & 110 & $P(11^2)$ \\
127 & 7  & 127 & 7   & Mersenne $2^7 - 1$ \\
354 & 58 & 59  & 58  & Full cycle \\
407 & 180 & 407 & 180 & $P(11 \times 37)$ \\
\bottomrule
\end{tabular}
\end{table}

\subsection{The 354-Anatomy and CRT Intertwining}\label{sec:354}

The integer $354 = 2 \times 3 \times 59$ plays a distinctive role in
the OMARWA arithmetic. Since $\gcd(354, 6) = 6$, the Core Reduction
gives:
\[
  P(354) = \ord_{59}(2) = 58.
\]

\begin{proposition}[354-Decomposition]\label{prop:354}
$354 = 6 \times 59 = 2 \times 177 = 2 \times 3 \times 59$, with
$177 = 3 \times 59$ acting as the antipodal modulus.
\end{proposition}

\begin{definition}[Intertwining Map]\label{def:intertwine}
The \emph{CRT intertwining map} is
$\iota : \mathbb{Z}_{354} \to \mathbb{Z}_{354}$ defined by
$\iota(s) = (60s + 295) \bmod 354$.
\end{definition}

This map arises from the CRT isomorphism
$\mathbb{Z}_{354} \cong \mathbb{Z}_6 \times \mathbb{Z}_{59}$ and
intertwines the mod~$6$ and mod~$59$ components. The orbit length of
the OMARWA residues modulo~$354$ equals~$58$, confirming
$P(354) = 58 = \ord_{59}(2)$.

\begin{proposition}[Ribbon Synchronization]\label{prop:ribbon-sync}
$\lcm(P(121), P(354)) = \lcm(110, 58) = 3190 = 110 \times 29$.
The global synchronization period is
$\lcm(L', P(354)) = \lcm(330, 58) = 9570 = 330 \times 29$.
\end{proposition}

% ══════════════════════════════════════════════════════════════
\section{Centered Hexagonal Intersections}
% ══════════════════════════════════════════════════════════════

The centered hexagonal numbers are $H_n = 3n(n-1) + 1$ for $n \geq 0$
(OEIS A003215~\cite{hex}).

\begin{proposition}\label{prop:h2}
$T_0 = H_2 = 7$.
\end{proposition}

\begin{proposition}\label{prop:h6}
$T_0 \cdot T_1 = 7 \times 13 = 91 = H_6$.
\end{proposition}

\begin{proposition}\label{prop:h4}
$T_3 = 49 = 7^2 = T_0^2$ and $T_3 = H_4$.
\end{proposition}

\begin{proposition}[Hexagonal Confinement]\label{prop:hex-conf}
$H_n \equiv 1 \pmod{6}$ for all $n \geq 0$, so every centered
hexagonal number lies on the $6n{+}1$ skeleton.
\end{proposition}

\begin{conjecture}[Finiteness]\label{conj:finite}
$\{k : T_k = H_n \text{ for some } n\} = \{0, 3\}$.
\end{conjecture}

% ══════════════════════════════════════════════════════════════
\section{Digital Root Structure}
% ══════════════════════════════════════════════════════════════

\begin{proposition}\label{prop:dr}
For all $k \geq 0$:
\[
  T_k \bmod 9 = \begin{cases} 7 & k \text{ even}, \\ 4 & k \text{ odd}.
  \end{cases}
\]
\end{proposition}

\begin{proof}
$P(9) = 2$: $T_0 \bmod 9 = 7$, $T_1 \bmod 9 = 4$, repeating.
\end{proof}

\begin{corollary}\label{cor:dr}
The digital root of $T_k$ alternates between $7$ (even~$k$) and $4$
(odd~$k$), with $4 + 7 = 11$.
\end{corollary}

% ══════════════════════════════════════════════════════════════
\section{Pascal--Sierpi\'nski--Fibonacci Bridge and the Collatz
  Connection}\label{sec:pascal}
% ══════════════════════════════════════════════════════════════

\subsection{Pascal--Sierpi\'nski Bridge}

\begin{proposition}[Carry Criterion]\label{prop:kummer}
The $p$-adic valuation $\nu_p\binom{T_i}{T_j}$ equals the number of
carries when adding $T_j$ and $T_i - T_j$ in base~$p$
(Kummer's theorem~\cite{kummer}).
\end{proposition}

\begin{proposition}[Sierpi\'nski Gasket]\label{prop:sierpinski}
$\binom{T_k}{j} \bmod 2$ produces Sierpi\'nski-type fractal patterns
at rows indexed by OMARWA terms.
\end{proposition}

\begin{theorem}[Fibonacci Residues in Mod 11]\label{thm:fib-mod11}
$T_1 \bmod 11 = F_3 = 2$,\;
$T_2 \bmod 11 = F_4 = 3$,\;
$T_3 \bmod 11 = F_5 = 5$.
\end{theorem}

\begin{proposition}[Lucas--Golden Bridge]\label{prop:lucas}
$L_5 = 11$ (Lucas number), $P(11) = 10 = \varphi(11)$, and
$F_{10} = F_5 \cdot L_5 = 5 \times 11 = 55$.
\end{proposition}

\begin{proposition}[CRT Pascal Decomposition]\label{prop:crt-pascal}
$\binom{T_k}{j} \bmod 6$ decomposes via CRT into mod~$2$ and mod~$3$
components, yielding Sierpi\'nski patterns respecting P622
crystallographic symmetry.
\end{proposition}

\subsection{Collatz--Fibonacci Bridge}\label{sec:collatz}

\begin{theorem}[Collatz--OMARWA Bridge]\label{thm:collatz}
The Collatz trajectory of $T_0 = 27$ (noting $27 = 3^3$) requires
exactly $111 = 3 \times 37$ steps to reach~$1$. The trajectory of
$T_4 = 97$ requires $118 = 2 \times 59$ steps. The difference is
$118 - 111 = 7 = T_0$.
\end{theorem}

\begin{theorem}[Fibonacci Order]\label{thm:fib-order}
$\ord_{577}(2) = 144 = F_{12}$, where $577$ is prime. This connects
multiplicative orders to the Fibonacci sequence through the Pisano
period.
\end{theorem}

% ══════════════════════════════════════════════════════════════
\section{Palindromic Symmetry and Fractal Holon
  Structures}\label{sec:palindrome}
% ══════════════════════════════════════════════════════════════

\subsection{Palindromic Involution}

\begin{theorem}[Palindromic Involution]\label{thm:palindrome}
For any $m$ with $P(m)$ even, the residue sequence
$r_k = T_k \bmod m$ satisfies
$r_k + r_{k + P(m)/2} \equiv 2 \pmod{m'}$ for all~$k$, where
$m' = m/\gcd(m,6)$.
\end{theorem}

\begin{proof}
Since $2^{P(m)} \equiv 1 \pmod{m'}$ and $P(m)$ is even,
$2^{P(m)/2} \equiv -1 \pmod{m'}$. Thus
$r_{k+P(m)/2} \equiv -6 \cdot 2^k + 1 \pmod{m'}$, giving
$r_k + r_{k+P(m)/2} \equiv 2 \pmod{m'}$.
\end{proof}

\begin{proposition}[29-1-29 Ribbon]\label{prop:ribbon29}
The full OMARWA orbit modulo~$59$ has period~$58$, decomposing as
$29 + 1 + 29 - 1 = 58$: two wings of length~$29$ separated by a
center and an antipodal node, where $29 = h(E_8) - 1$ is the maximal
Weyl exponent of~$E_8$.
\end{proposition}

\begin{proposition}[354--453 Cosmic Palindrome]\label{prop:354-453}
$354 \equiv 453 \pmod{9}$ and $354 \equiv 453 \pmod{11}$, with
$453 - 354 = 99 = 9 \times 11$.
The digit-reversal $354 \leftrightarrow 453$ is a palindromic
involution that preserves both the mod~$9$ and mod~$11$ residue classes.
\end{proposition}

\begin{proposition}[Emirp Palindrome]\label{prop:emirp-palindrome}
$37$ and $73$ are an emirp pair. The palindromic prime $313$ divides
$37 \times 73 - 2388 = 313$, and $313$ is indeed a palindromic prime.
\end{proposition}

\subsection{Fractal Holon Structures}

\begin{theorem}[Fractal Period Conservation]\label{thm:fractal-holon}
$P(111) = P(37) = 36$ and $P(333) = 36$. The period is conserved under
the Trinity multiplication $37 \to 3 \times 37 = 111 \to
3 \times 111 = 333$ because $\gcd(3, 37) = 1$ and $P(3) = 1$.
\end{theorem}

\begin{proposition}[729-Period]\label{prop:729}
$P(729) = P(3^6) = 2 \cdot 3^4 = 162$.
\end{proposition}

\begin{proposition}[Quadratic Residue]\label{prop:quad-res}
$\ord_{11}(4) = 5$, governing the mod~$11$ behavior of $37^n$
since $37 \equiv 4 \pmod{11}$.
\end{proposition}

% ══════════════════════════════════════════════════════════════
\section{Network Topology and Holon
  Crystallography}\label{sec:holon}
% ══════════════════════════════════════════════════════════════

\subsection{354-Plane Indexing and Network Theorems}

The OMARWA arithmetic generates a family of network topology identities
(MT-1 through MT-14) based on the $354$-plane indexing:

\begin{theorem}[354-Plane Identity (MT-1)]\label{thm:mt1}
$354 = 6 \times 59$, with each ``sector'' of~$59$ planes decomposing as
$59 = 29 + 1 + 29$ (ribbon half-width + center + ribbon half-width).
\end{theorem}

\begin{theorem}[Antipodal Involution (MT-3)]\label{thm:mt3}
The map $i \mapsto (i + 177) \bmod 354$ is an involution on
$\mathbb{Z}_{354}$, partitioning the $354$ planes into $177$ antipodal
pairs.
\end{theorem}

\begin{theorem}[29-1-29 Ribbon Decomposition (MT-4)]\label{thm:mt4}
Each sector contains $59$ planes: $29$ in the left wing, $1$ central
junction, and $29$ in the right wing. The value $29 = h(E_8) - 1$
is the maximal Weyl exponent of~$E_8$.
\end{theorem}

\subsection{Holon Crystallography}

\begin{theorem}[Trinity Palindrome]\label{thm:trinity}
$37 + 111 + 111 + 111 + 37 = 407 = 11 \times 37$.
\end{theorem}

\begin{proof}
Direct computation: $37 + 3 \times 111 + 37 = 74 + 333 = 407$, and
$407 / 11 = 37$.
\end{proof}

\begin{proposition}[Holon Hierarchy]\label{prop:holon-hier}
$37 = H_4$ (centered hexagonal), $111 = 3 \times 37$,
$333 = 9 \times 37$, $370 = 10 \times 37$, and $370 \bmod 11 = 7 = T_0$
(closure).
\end{proposition}

\begin{theorem}[Double Lock]\label{thm:double-lock}
$T_k \equiv 8 \pmod{11}$ implies $T_k \equiv 1 \pmod{8}$.
This ``APEX triggering'' occurs at $k \equiv 8 \pmod{10}$; for
$k \geq 3$, $T_k \bmod 8 = 1$ always, since $2^k \bmod 8 = 0$.
\end{theorem}

\begin{proposition}[Pisano Synchronization]\label{prop:pisano}
The Pisano period $\pi(11) = 10 = P(11)$: the Fibonacci sequence
modulo~$11$ and the OMARWA sequence modulo~$11$ share the same
fundamental period.
\end{proposition}

% ══════════════════════════════════════════════════════════════
\section{Weyl Exponents of
  \texorpdfstring{$E_8$}{E8}}\label{sec:weyl}
% ══════════════════════════════════════════════════════════════

The exponents of the Weyl group of~$E_8$ are
$\{m_1, \ldots, m_8\} = \{1, 7, 11, 13, 17, 19, 23, 29\}$.

\begin{theorem}[Weyl Exponent Sum]\label{thm:weyl-sum}
$\displaystyle\sum_{i=1}^{8} m_i = 1 + 7 + 11 + 13 + 17 + 19 + 23 + 29
= 120 = \frac{\dim(E_8)}{2} - 4$.
\end{theorem}

\begin{theorem}[Palindromic Pairing]\label{thm:weyl-palindrome}
$m_i + m_{9-i} = 30 = h(E_8)$ for $i = 1, \ldots, 4$:
\[
  1 + 29 = 7 + 23 = 11 + 19 = 13 + 17 = 30.
\]
\end{theorem}

\begin{proposition}[Coprimality]\label{prop:weyl-coprime}
Every Weyl exponent of~$E_8$ is coprime to $30 = h(E_8)$; equivalently,
the exponents are exactly the units $(\mathbb{Z}/30\mathbb{Z})^{\times}$
with $\varphi(30) = 8 = \mathrm{rank}(E_8)$.
\end{proposition}

\begin{theorem}[Dimension Formula]\label{thm:dim-e8}
$\dim(E_8) = 248 = 2 \times 120 + 8 = 2\sum m_i + \mathrm{rank}(E_8)$.
\end{theorem}

\begin{remark}
The OMARWA connection: $m_2 = 7 = T_0$ and $m_8 = 29 = h(E_8) - 1$.
The palindromic pairing $m_i + m_{9-i} = h(E_8)$ mirrors the
palindromic involution of OMARWA residues
(Theorem~\ref{thm:palindrome}).
\end{remark}

% ══════════════════════════════════════════════════════════════
\section{Exceptional Universality and Coxeter
  Numbers}\label{sec:exceptional}
% ══════════════════════════════════════════════════════════════

\begin{theorem}[Coxeter--Period Identity]\label{thm:coxeter}
$h(E_8) = 30 = L$ (the OMARWA super-period).
Moreover, $\varphi(h(E_8)) = \varphi(30) = 8 = \mathrm{rank}(E_8)$.
\end{theorem}

\begin{theorem}[Exceptional Universality]\label{thm:exceptional}
The Coxeter numbers of all five exceptional simple Lie groups arise from
OMARWA arithmetic:
\end{theorem}

\begin{table}[H]
\centering
\caption{Exceptional Coxeter numbers from OMARWA periods.}
\label{tab:exceptional}
\smallskip
\begin{tabular}{@{}lccl@{}}
\toprule
Group $G$ & $h(G)$ & OMARWA source & Derivation \\
\midrule
$G_2$ & 6  & $P(27) = 6$     & $\ord_9(2) = 6$ \\
$F_4$ & 12 & $|D_6| = 12$    & $2 \times P(27) = 2 \times h(G_2)$ \\
$E_6$ & 12 & $|D_6| = 12$    & P622 point group order \\
$E_7$ & 18 & $P(81) = 18$    & $\ord_{27}(2) = 18$ \\
$E_8$ & 30 & $L = P(77) = 30$ & $\lcm(P(7), P(11))$ (super-period) \\
\bottomrule
\end{tabular}
\end{table}

\begin{proof}
Each Coxeter number~$h(G)$ matches an OMARWA period:
\begin{itemize}
  \item $P(27) = P(3^3) = \ord_9(2) = 6 = h(G_2)$;
  \item $P(81) = P(3^4) = \ord_{27}(2) = 18 = h(E_7)$;
  \item $L = P(77) = 30 = h(E_8)$.
\end{itemize}
For $F_4$ and $E_6$ (both $h = 12$): $12 = 2 \times h(G_2)$ arises as
the order of the diheral group~$D_6$, the point group of the P622
hexagonal lattice. That $h(G_2) = 6$, $h(E_7) = 18$, $h(E_8) = 30$ are
standard~\cite{humphreys}. The Weyl group order
$|W(E_8)| = 696{,}729{,}600$ and the Rank Lock $\varphi(30) = 8$
provide additional structural confirmation.
\end{proof}

% ══════════════════════════════════════════════════════════════
\section{Topological and Physical
  Interpretations}\label{sec:topology}
% ══════════════════════════════════════════════════════════════

The algebraic structures established in Sections~2--11 admit natural
topological and physical interpretations. We present these as
\emph{model statements} (epistemic class~[B]/[C]) rather than proven
theorems; all underlying algebraic identities are verified in Lean~4.

\subsection{Winding Numbers on \texorpdfstring{$T^3$}{T3}}

\begin{proposition}[Toroidal Winding]\label{prop:winding}
The super-period $L = 30$ traces a path on the $3$-torus
$T^3 = S^1 \times S^1 \times S^1$ where each circle corresponds to one
modular period:
\[
  (n_1, n_2, n_3) = \left(\frac{L}{P(9)},\;
  \frac{L}{P(11)},\; \frac{L}{P(27)}\right)
  = (15, 3, 5).
\]
The winding number product is $15 \times 3 \times 5 = 225 = 15^2$, and
the sum $15 + 3 + 5 = 23$ is prime.
\end{proposition}

\subsection{Chern--Simons Level}

\begin{proposition}[Chern--Simons Level]\label{prop:cs}
Setting the Chern--Simons level $k_{\mathrm{CS}} = 30 = L = h(E_8)$
yields:
\begin{itemize}
  \item Berry phase: $\gamma_{30} = 15\pi$ (half-integer winding),
  \item First Chern number: $C_1 = k_{\mathrm{CS}} / 2 = 15$,
  \item Anyonic exchange angle: $\theta = \pi / k_{\mathrm{CS}}
    = \pi / 30$,
  \item Ground state degeneracy on $T^2$: $\mathrm{GSD} = k_{\mathrm{CS}} = 30$.
\end{itemize}
\end{proposition}

The three independent derivations of $k_{\mathrm{CS}} = 30$ (from the
Coxeter number, from the super-period, and from the prime
factorization $30 = 2 \times 3 \times 5$) are verified algebraically in
Lean~4.

\subsection{Hopf Fibration Structure}

\begin{proposition}[Hopf Reduction]\label{prop:hopf}
The antipodal modulus $177 = 354/2 = 3 \times 59$ defines a Hopf-type
fibration $\Sigma^3 \to S^2_{\mathrm{OMARWA}}$ with:
\begin{itemize}
  \item fiber identification via the involution
    $i \mapsto (i + 177) \bmod 354$,
  \item ribbon width $59 = 2h(E_8) - 1$,
  \item half-width $29 = h(E_8) - 1 = m_8$ (maximal Weyl exponent).
\end{itemize}
\end{proposition}

\subsection{Physical Scaling}

\begin{proposition}[APEX Triggering]\label{prop:apex}
$T_k \equiv 8 \pmod{11}$ occurs exactly $3$ times per super-period
($k \equiv 8 \pmod{10}$ within $[0, 29]$: at $k = 8, 18, 28$).
These ``APEX positions'' mark the mod~$11$ residue closest to the
forbidden unity.
\end{proposition}

\begin{proposition}[Fractional Order Convergence]\label{prop:alpha}
The ratio $\alpha_k = T_k / T_{k+1}$ converges to $1/2$ as
$k \to \infty$, with
$\alpha_k = (6 \cdot 2^k + 1) / (6 \cdot 2^{k+1} + 1) \to 1/2$.
\end{proposition}

% ══════════════════════════════════════════════════════════════
\section{Conclusion and Open Questions}
% ══════════════════════════════════════════════════════════════

We have established a complete modular periodicity theory for the
sequence $T_k = 6 \cdot 2^k + 1$, reducing all period computations to
multiplicative orders via the Core Reduction Theorem. The fractal period
laws, product theorem, and super-period $L = 30$ reveal an unexpectedly
rich arithmetic structure. The Forbidden Unity Theorem, CRT
intertwining, palindromic involutions, and network topology identities
demonstrate that this elementary sequence encodes deep number-theoretic
relationships. The coincidence of all five exceptional Coxeter numbers
with OMARWA-derived invariants (Exceptional Universality) is a striking
observation whose structural explanation, if any, remains open.

All principal results are formally verified in Lean~4 with zero
\texttt{sorry} axioms across 28 modules containing 956 theorems.

\subsection*{Open Questions}

\begin{enumerate}
  \item \textbf{Wieferich exception:} Does the fractal law fail for
    $p = 1093$ and $p = 3511$?

  \item \textbf{Finite intersections:} Is
    $\{k : T_k = H_n \text{ for some } n\} = \{0, 3\}$?

  \item \textbf{Structural interpretation:} Is there a
    representation-theoretic explanation for Exceptional Universality
    (Theorem~\ref{thm:exceptional})?

  \item \textbf{Artin density:} Does the density of primes $p$ with
    $P(p) = p - 1$ among OMARWA-related moduli differ from the Artin
    constant?

  \item \textbf{Generalized sequences:} For which $(a,b) \in
    \mathbb{Z}^2$ does $S_k(a,b) = a \cdot 2^k + b$ have super-period
    equal to a Coxeter number?

  \item \textbf{Forbidden unity radius:} How does the ``distance'' from
    the forbidden residue~$1$ in the mod~$m$ orbit behave as
    $m \to \infty$?

  \item \textbf{Collatz bridge:} Is the coincidence
    $\mathrm{Collatz}(27) = 111 = 3 \times 37$ a systematic feature of
    Proth numbers, or an isolated curiosity?

  \item \textbf{Topological content:} Does the toroidal winding
    $(15, 3, 5)$ have a physical realization beyond the formal
    $T^3$~model?
\end{enumerate}

% ══════════════════════════════════════════════════════════════
\appendix
% ══════════════════════════════════════════════════════════════

\section{Lean 4 Proof Catalog}\label{app:lean}

All theorems in this paper have been formally verified using the
\textbf{OmarwaCore} library:

\medskip
\begin{center}
\begin{tabular}{@{}ll@{}}
\toprule
\textbf{Engine}       & Lean 4.29.0-rc3 with Mathlib4 \\
\textbf{Theorems + Lemmas} & 956 \\
\textbf{Sorry axioms} & 0 (zero) \\
\textbf{Source lines}  & 5,779 \\
\textbf{Modules}       & 28 \\
\bottomrule
\end{tabular}
\end{center}

\smallskip
\noindent Full source:
\href{https://github.com/omarwa622/OmarwaCore}{github.com/omarwa622/OmarwaCore}

\smallskip
\noindent Complete module-to-section mapping:

\begin{table}[H]
\centering
\small
\caption{All 28 Lean~4 modules and their paper locations.}\label{tab:lean}
\smallskip
\begin{tabular}{@{}rlll@{}}
\toprule
\# & Lean Module & Section & Key Result \\
\midrule
1  & \texttt{Sequence}             & \S2.1  & $T_k$ definition, recurrence \\
2  & \texttt{ModularPeriods}       & \S2.2  & $P(m)$ values, mod 11 \\
3  & \texttt{SabbathNode}          & \S2.3  & $T_6 = 5 \times 7 \times 11$ \\
4  & \texttt{FractalLaws}          & \S3.1  & Prime-power periods \\
5  & \texttt{SuperPeriod}          & \S3.5  & $L = 30$ \\
6  & \texttt{CompositeModular}     & \S3.6  & Composite $P(m)$ \\
7  & \texttt{MersenneFermat}       & \S3.3  & $P(7)=3, P(31)=5, P(17)=8$ \\
8  & \texttt{ForbiddenUnityRibbon} & \S3.4  & $T_k \not\equiv 1$; ribbon \\
9  & \texttt{AffineOperator}       & \S4    & $6n+1$ skeleton, sampling \\
10 & \texttt{ProductTheorem}       & \S5.1  & CRT product theorem \\
11 & \texttt{DeepArithmetic}       & \S5.1  & GCD law, $\varphi$ values \\
12 & \texttt{CRTIntertwining}      & \S5.2  & 354-anatomy, $\iota$ map \\
13 & \texttt{CenteredHexagonal}    & \S6    & $H_2 = T_0$, $H_6 = T_0 T_1$ \\
14 & \texttt{ResidueDR}            & \S7    & Digital root bisection \\
15 & \texttt{PascalSierpinski}     & \S8.1  & Kummer, Sierpi\'nski gasket \\
16 & \texttt{GrandSynthesis}       & \S8.1  & Fibonacci mod 11 embedding \\
17 & \texttt{CollatzFibonacci}     & \S8.2  & Collatz(27)=111, $F_{12}$ \\
18 & \texttt{PalindromicInvolution} & \S9.1 & 354$\leftrightarrow$453, emirps \\
19 & \texttt{FractalHolon}         & \S9.2  & $P(111)=P(333)=36$ \\
20 & \texttt{NetworkTheorems}      & \S10.1 & MT-1\ldots MT-14, 354-plane \\
21 & \texttt{HolonCrystallographic} & \S10.2 & Trinity $407=11\times37$ \\
22 & \texttt{WeylExponents}        & \S11   & $\sum m_i = 120$, palindrome \\
23 & \texttt{CoxeterNumber}        & \S12   & $h(E_8)=L=30$, rank lock \\
24 & \texttt{ExceptionalUniversality} & \S12 & All 5 exceptionals \\
25 & \texttt{ChernSimonsLevel}     & \S12.2 & $k_{\mathrm{CS}}=30$ \\
26 & \texttt{HopfFibration}        & \S12.3 & Antipodal 177, ribbon 59 \\
27 & \texttt{TopologicalWinding}   & \S12.1 & $(15,3,5)$ on $T^3$ \\
28 & \texttt{PhysicalScaling}      & \S12.4 & APEX, $\alpha_k \to 1/2$ \\
\bottomrule
\end{tabular}
\end{table}

% ══════════════════════════════════════════════════════════════
\section{Computational Verification}\label{app:comp}

All results have been independently verified computationally:

\begin{itemize}
  \item \textbf{Period table:} $P(m)$ for $m = 1, \ldots, 10{,}000$
  \item \textbf{GCD matrix:} $\gcd(T_i, T_j)$ for
    $0 \leq i, j \leq 100$
  \item \textbf{Fractal law:} $P(p^n) = \ord_p(2) \cdot p^{n-1}$
    for all primes $p < 1000$ (excluding $1093$) and $n \leq 5$
  \item \textbf{Digital root:} $\DR(T_k) \in \{4, 7\}$ verified for
    $k \leq 10{,}000$
  \item \textbf{Palindromic symmetry:} All $m \leq 1000$ with
    even $P(m)$
  \item \textbf{Forbidden unity:} $T_k \bmod m \neq 1$ verified for all
    $m \leq 1000$ with $\gcd(m,6) = 1$ and $k \leq P(m)$
  \item \textbf{Network identities:} MT-1 through MT-14 verified
    numerically
  \item \textbf{Collatz trajectories:} Collatz($27$) = 111 steps,
    Collatz($97$) = 118 steps independently verified
  \item \textbf{Weyl exponents:} All eight palindromic pairings
    $m_i + m_{9-i} = 30$ verified, $\sum m_i = 120$
  \item \textbf{354--453 palindrome:} Congruences mod $9$ and mod $11$
    verified
\end{itemize}

\noindent Dataset and scripts:
\href{https://github.com/omarwa622/OmarwaCore}{github.com/omarwa622/OmarwaCore}.

% ══════════════════════════════════════════════════════════════
% REFERENCES
% ══════════════════════════════════════════════════════════════

\begin{thebibliography}{14}

\bibitem{oeis}
OEIS Foundation,
``A004119: $a(0) = 1$; thereafter $a(n) = 3 \cdot 2^{n-1} + 1$,''
\textit{The On-Line Encyclopedia of Integer Sequences},
\url{https://oeis.org/A004119}.

\bibitem{artin}
E.~Artin,
\textit{Collected Papers},
Springer, 1965.

\bibitem{ireland}
K.~Ireland and M.~Rosen,
\textit{A Classical Introduction to Modern Number Theory}, 2nd~ed.,
Springer GTM~84, 1990.

\bibitem{humphreys}
J.\,E.~Humphreys,
\textit{Reflection Groups and Coxeter Groups},
Cambridge Studies in Advanced Mathematics~29, 1990.

\bibitem{mathlib}
The Lean Community,
``Mathlib4: The Mathematics Library for Lean~4,''
\url{https://github.com/leanprover-community/mathlib4}.

\bibitem{omarwacore}
\"O.~\c{C}etinta\c{s},
``OmarwaCore: Lean~4 Formalization of the OMARWA Sequence,''
2026.
\url{https://github.com/omarwa622/OmarwaCore}.

\bibitem{proth}
F.~Proth,
``Th\'eor\`emes sur les nombres premiers,''
\textit{Comptes Rendus Acad.\ Sci.\ Paris},
vol.~87, pp.~926--928, 1878.

\bibitem{narkiewicz}
W.~Narkiewicz,
\textit{Elementary and Analytic Theory of Algebraic Numbers},
3rd~ed., Springer, 2004.

\bibitem{wieferich}
F.\,G.~Dorais and D.~Klyve,
``A Wieferich prime search up to $6.7 \times 10^{15}$,''
\textit{J.\ Integer Seq.},
vol.~14, Article~11.9.2, 2011.

\bibitem{hex}
OEIS Foundation,
``A003215: Hex (or centered hexagonal) numbers,''
\textit{The On-Line Encyclopedia of Integer Sequences},
\url{https://oeis.org/A003215}.

\bibitem{kummer}
E.\,E.~Kummer,
``\"Uber die Erg\"anzungss\"atze zu den allgemeinen
Reciprocit\"atsgesetzen,''
\textit{J.~reine angew.~Math.}, vol.~44, pp.~93--146, 1852.

\bibitem{golomb}
S.\,W.~Golomb,
``Properties of the sequence $3 \cdot 2^n + 1$,''
\textit{Math.\ Comp.},
vol.~30, no.~135, pp.~657--663, 1976.

\bibitem{bourbaki}
N.~Bourbaki,
\textit{Lie Groups and Lie Algebras: Chapters 4--6},
Springer, 2002.

\bibitem{nakahara}
M.~Nakahara,
\textit{Geometry, Topology and Physics}, 2nd~ed.,
IOP Publishing, 2003.

\end{thebibliography}

\bigskip
\noindent\rule{\textwidth}{0.4pt}

\smallskip
{\small\noindent\textit{Manuscript prepared: March 2026.
Formal verification: Lean~4.29.0-rc3 + Mathlib4 (956~theorems,
28~modules, 0~sorry).}\\
\textit{Source:
\href{https://github.com/omarwa622/OmarwaCore}{github.com/omarwa622/OmarwaCore}.}}

\end{document}
